\section{Anbindung externer Datenquellen}
Zur Erweiterung der Funktionalität werden externe Datenquellen eingebunden.  
Die Kommunikation erfolgt über eine \gls{wlan}-Verbindung.  
Der Verbindungsaufbau wird durch eine eigene Funktion realisiert:  
\begin{verbatim}
void connectToWiFi()
<HIER KOMMT NOCH EIN CODEAUSSCHNITT HIN>
\end{verbatim}
Die Zeit wird über \gls{ntp} synchronisiert und die Wetterdaten von der Web-\gls{api} von Open-Meteo.
Die Synchronisation erfolgt initial beim Systemstart sowie zyklisch im Betrieb.   
Die Daten werden über \gls{https} abgerufen und im \gls{json}-Format verarbeitet.  
Die zentrale Funktion zur Aktualisierung der Wetterdaten lautet:  
\begin{verbatim}
bool updateWeatherFromServer()
<HIER KOMMT NOCH EIN CODEAUSSCHNITT HIN>
\end{verbatim}
Die empfangenen Daten werden anschließend klassifiziert.  
Hierzu werden die Temperaturwerte über eine Case-Anweisung ausgewählt und der Wettercode in eine definierte Wettergruppe überführt.
Mit dieser wird dann ebenfalls über eine Case-Anweisung die aktuelle Wetteranzeige ausgewählt.
Die Anzeige erfolgt dann durch Zuordnung zu entsprechenden Wörtern welche als Globale Variablen deklariert sind.